\chapter{Emulasi dan PID Diskrit}
\label{bab:emulasi}
\enchap{Emulation and Discrete PID Control}

\begin{tujuan}
	Setelah mempelajari bab ini, Anda mampu mendiskritisasi pengontrol kontinu dengan berbagai metode, menjelaskan distorsi frekuensi Tustin dan \textit{prewarping}, menurunkan PID diskrit bentuk posisi dan kecepatan, menerapkan \textit{anti-windup} dan filter derivatif, serta menala PID diskrit secara heuristik.
\end{tujuan}

\section{Perancangan dengan emulasi}
Emulasi berarti merancang pengontrol kontinu $C(s)$ dengan metode yang sudah dikenal, lalu mengubahnya menjadi pengontrol diskrit $D(z)$ dan memverifikasi hasilnya dengan model diskrit plant. Kelebihannya, seluruh pengalaman kontrol kontinu dapat dipakai. Kekurangannya, tunda ZOH $T/2$ diabaikan saat merancang $C(s)$, sehingga hasilnya hanya baik bila $T$ kecil ($\omega_s \gtrsim 20\omega_b$).

\section{Metode diskritisasi}
\subsection{Aproksimasi integrasi}
Integrator $1/s$ dapat diganti dengan aturan integrasi numerik, yang setara dengan substitusi $s$ sebagai fungsi $z$:
\begin{align}
	\text{Euler maju:}\quad & s \leftarrow \frac{z-1}{T}, \\
	\text{Euler mundur:}\quad & s \leftarrow \frac{z-1}{zT}, \\
	\text{Tustin (bilinear):}\quad & s \leftarrow \frac{2}{T}\,\frac{z-1}{z+1}.
\end{align}
Ketiganya memetakan setengah bidang kiri bidang $s$ secara berbeda. Euler maju memetakannya ke setengah bidang $\text{Re}\,z < 1$, sehingga $C(s)$ yang stabil dapat menjadi $D(z)$ yang tidak stabil (misalnya $s = -50$ dengan $T = 0{,}05$\,s dipetakan ke $z = -1{,}5$). Euler mundur memetakannya ke lingkaran berpusat $0{,}5$ berjari-jari $0{,}5$, selalu stabil tetapi cenderung terlalu teredam. Tustin memetakan setengah bidang kiri tepat ke dalam lingkaran satuan, sehingga paling banyak dipakai.

\subsection{Distorsi frekuensi dan \textit{prewarping}}
Substitusi $z = e^{j\omega_dT}$ ke Tustin menghasilkan hubungan frekuensi
\begin{equation}
	\omega_c = \frac{2}{T}\tan\frac{\omega_dT}{2}.
\end{equation}
Seluruh sumbu $j\omega_c$ dimampatkan ke $0 \le \omega_d < \omega_N$. Distorsi kecil bila $\omega_d \ll \omega_N$. Agar satu frekuensi kritis $\omega_1$ (misalnya \textit{crossover} atau frekuensi \textit{notch}) dipetakan tepat, gunakan \textit{prewarping}:
\begin{equation}
	s \leftarrow \frac{\omega_1}{\tan(\omega_1T/2)}\,\frac{z-1}{z+1}.
\end{equation}

\subsection{\textit{Matched pole-zero}}
Metode MPZ memetakan setiap kutub dan zero berhingga dengan $z = e^{sT}$, memetakan zero di tak hingga ke $z = -1$ (atau membiarkan satu tunda agar kausal-tegas), lalu menyamakan penguatan DC $D(1) = C(0)$.

\begin{contoh}[Diskritisasi kompensator \textit{lead}]
	$C(s) = 10(s+2)/(s+10)$ dengan $T = 0{,}05$\,s menghasilkan:
	\begin{center}
		\renewcommand{\arraystretch}{2.0}
		\begin{tabular}{@{}l l@{}}
			\toprule
			Metode & $D(z)$ \\
			\midrule
			Euler maju & $\dfrac{10z - 9}{z - 0{,}5}$ \\
			Euler mundur & $\dfrac{7{,}333z - 6{,}667}{z - 0{,}6667}$ \\
			Tustin & $\dfrac{8{,}4z - 7{,}6}{z - 0{,}6}$ \\
			MPZ & $8{,}269\,\dfrac{z - 0{,}9048}{z - 0{,}6065}$ \\
			\bottomrule
		\end{tabular}
	\end{center}
	Pada MPZ, zero $e^{-2T} = 0{,}9048$, kutub $e^{-10T} = 0{,}6065$, dan penguatan dari $D(1) = C(0) = 2$. Di bawah sekitar 5\,rad/s respons frekuensi keempatnya hampir sama; mendekati $\omega_N = 62{,}8$\,rad/s perbedaannya besar, terutama pada fasa. Tustin dan MPZ paling dekat dengan $C(s)$ di rentang kerja.
\end{contoh}

\section{PID diskrit}
\subsection{Bentuk posisi}
Dengan integral persegi mundur dan derivatif beda mundur, PID menjadi
\begin{equation}
	u[k] = K_p\Big[e[k] + \frac{T}{T_i}\sum_{j=0}^{k}e[j] + \frac{T_d}{T}\big(e[k]-e[k-1]\big)\Big],
\end{equation}
yang diimplementasikan secara rekursif dengan $I[k] = I[k-1] + \frac{K_pT}{T_i}e[k]$. Fungsi transfernya
\begin{equation}
	D(z) = K_p\left[1 + \frac{T}{T_i}\,\frac{z}{z-1} + \frac{T_d}{T}\,\frac{z-1}{z}\right].
\end{equation}
Integrator diskrit adalah kutub di $z = 1$. Perhatikan bahwa penguatan derivatif efektif $K_pT_d/T$ membesar bila $T$ kecil, sehingga derau ikut membesar.

\subsection{Bentuk kecepatan}
Mengurangkan $u[k-1]$ dari $u[k]$ menghasilkan bentuk inkremental
\begin{equation}
	\Delta u[k] = q_0e[k] + q_1e[k-1] + q_2e[k-2],
\end{equation}
\begin{equation}
	q_0 = K_p\Big(1 + \frac{T}{T_i} + \frac{T_d}{T}\Big),\qquad q_1 = -K_p\Big(1 + \frac{2T_d}{T}\Big),\qquad q_2 = K_p\frac{T_d}{T},
\end{equation}
dengan $u[k] = u[k-1] + \Delta u[k]$. Contoh: $K_p = 2$, $T_i = 0{,}5$\,s, $T_d = 0{,}05$\,s, $T = 0{,}1$\,s memberi $q_0 = 3{,}4$, $q_1 = -4{,}0$, $q_2 = 1{,}0$. Bentuk kecepatan tidak memiliki akumulator eksplisit, memudahkan transfer manual--otomatis tanpa lonjakan, dan \textit{anti-windup} cukup dilakukan dengan membatasi $u[k]$.

\subsection{PID praktis}
Implementasi PID yang andal memerlukan beberapa modifikasi.
\begin{enumerate}
	\item \textbf{Filter derivatif}. $T_ds \rightarrow T_ds/(1 + T_ds/N)$; dengan Euler mundur dan $T_f = T_d/N$,
	\[
		D_f[k] = \frac{T_fD_f[k-1] + T_d\big(y[k]-y[k-1]\big)}{T_f + T}.
	\]
	\item \textbf{Derivatif pada pengukuran}: gunakan $-y$ alih-alih $e$ agar perubahan tangga \textit{setpoint} tidak menimbulkan lonjakan.
	\item \textbf{Bobot \textit{setpoint}}: bagian proporsional memakai $\beta r - y$ dengan $0 \le \beta \le 1$.
	\item \textbf{\textit{Anti-windup}}. Saat aktuator jenuh, integrator terus menumpuk galat sehingga menimbulkan lewatan besar. \textit{Clamping} menghentikan integrasi bila $u$ jenuh dan galat mendorong lebih jenuh; \textit{back-calculation} menambahkan $\frac{T}{T_t}(u_{\text{sat}} - u)$ ke integrator.
	\item \textbf{Batas nilai dan laju} $u$ sesuai aktuator (kit: $-100$ hingga $100\%$).
\end{enumerate}

\begin{contoh}[\textit{Integrator windup}]
	Plant kecepatan acuan dikontrol PI dengan $K_p = 3$, $T_i = 0{,}15$\,s, $T = 0{,}1$\,s, aktuator jenuh pada $|u| \le 1{,}2$, dan $r = 1$. Tanpa \textit{anti-windup}, $M_p \approx 19{,}6\%$ dan $t_s \approx 3{,}2$\,s. Dengan \textit{clamping}, $M_p \approx 0{,}3\%$ dan $t_s \approx 1{,}0$\,s (Gambar~\ref{fig:b5-windup}).
\end{contoh}

\begin{figure}[H]
	\centering
	\begin{tikzpicture}
		\begin{groupplot}[group style={group size=2 by 1,horizontal sep=1.5cm},skd,width=7cm,height=4.6cm,xmin=0,xmax=6,xlabel={$t$ (s)}]
			\nextgroupplot[ylabel={$y$},ymin=0,ymax=1.3,legend style={at={(0.98,0.05)},anchor=south east}]
			\addplot[skdgray,dashed,thin,domain=0:6,forget plot] {1};
			\addplot[skdred,const plot] table[x=t,y=yn] {\dat{p5_windup.dat}};
			\addplot[skdblue,const plot] table[x=t,y=yc] {\dat{p5_windup.dat}};
			\legend{tanpa AW, \textit{clamping}}
			\nextgroupplot[ylabel={$u$},ymin=0,ymax=1.4]
			\addplot[skdred,const plot] table[x=t,y=un] {\dat{p5_windup.dat}};
			\addplot[skdblue,const plot] table[x=t,y=uc] {\dat{p5_windup.dat}};
		\end{groupplot}
	\end{tikzpicture}
	\caption{Pengaruh \textit{anti-windup} pada PI dengan aktuator jenuh}
	\label{fig:b5-windup}
\end{figure}

\begin{lstlisting}[language=C++,caption={PID diskrit dengan D pada pengukuran dan \textit{clamping}},label={lst:b5-pid}]
float pid(float r, float y) {        // dipanggil tiap T
  float e = r - y;
  float P = Kp * e;
  Df = (Tf*Df + Td*(y - yPrev)) / (Tf + T);
  float D = -Kp * Df;                // D pada y
  float v = P + I + D;
  float u = constrain(v, -100, 100);
  if (u == v || e*v < 0)             // clamping
    I += Kp * T / Ti * e;
  yPrev = y;
  return u;
}
\end{lstlisting}

\section{Penalaan heuristik PID diskrit}
Karena ZOH menambah tunda sekitar $T/2$, tabel Ziegler--Nichols dan Cohen--Coon dipakai dengan tunda terkoreksi
\begin{equation}
	\theta' = \theta + \frac{T}{2}.
\end{equation}

\begin{contoh}[Koreksi tunda ZOH]
	Untuk plant FOPDT Bab~1 dengan $T = 0{,}05$\,s, penalaan dengan $\theta = 0{,}3$\,s ($K_p = 3{,}00$, $T_i = 0{,}60$\,s) menghasilkan $M_p \approx 89\%$ dan respons belum masuk pita 2\% dalam 8\,s. Dengan $\theta' = 0{,}325$\,s ($K_p = 2{,}77$, $T_i = 0{,}65$\,s, $T_d = 0{,}163$\,s) diperoleh $M_p \approx 72\%$ dan $t_s \approx 6{,}0$\,s. Koreksi membantu, tetapi ZN tetap agresif sehingga penalaan halus tetap diperlukan.
\end{contoh}

Prosedur penalaan halus: mulai dari hasil tabel dengan $\theta'$; turunkan $K_p$ hingga lewatan sesuai spesifikasi; atur $T_i$ (kecilkan untuk menghilangkan galat lebih cepat, besarkan bila berosilasi); tambahkan $T_d$ secukupnya sambil memperhatikan derau pada $u$; uji perubahan \textit{setpoint} naik, turun, dan gangguan beban; catat $M_p$, $t_s$, $e_{ss}$, IAE, dan $u_{\max}$. Prosedur ini dipraktikkan pada Modul 2.

\begin{ringkasan}
	Emulasi: rancang $C(s)$, lalu diskritisasi dengan Tustin atau MPZ; valid bila $T$ kecil. Euler maju dapat menghasilkan $D(z)$ tidak stabil. PID diskrit tersedia dalam bentuk posisi dan kecepatan; filter derivatif, derivatif pada pengukuran, dan \textit{anti-windup} wajib dalam praktik. Untuk penalaan heuristik, tambahkan $T/2$ pada tunda plant.\\[3pt]
	\en{Emulation discretises a continuous design (Tustin or matched pole-zero preferred). Discrete PID comes in position and velocity forms; practical PID needs derivative filtering, derivative on measurement, and anti-windup. Add $T/2$ to the process delay when using tuning tables.}
\end{ringkasan}

\paragraph{Istilah kunci.} emulasi \en{(emulation)}, transformasi bilinear (Tustin), \textit{prewarping}, \textit{matched pole-zero}, bentuk posisi/kecepatan \en{(position/velocity form)}, \textit{integrator windup}, \textit{anti-windup}, penalaan halus \en{(fine tuning)}.

\section*{Latihan}
\begin{enumerate}
	\item Diskritisasi $C(s) = 5(s+1)/(s+20)$ dengan Tustin dan MPZ untuk $T = 0{,}02$\,s, lalu bandingkan respons frekuensinya.
	\item Turunkan $q_0$, $q_1$, $q_2$ bentuk kecepatan dari bentuk posisi.
	\item Tunjukkan bahwa Euler maju memetakan $s = -50$ ke luar lingkaran satuan bila $T = 0{,}05$\,s.
	\item Turunkan rumus \textit{prewarping} dan terapkan untuk filter \textit{notch} 50\,Hz dengan $T = 1$\,ms.
	\item Simulasikan PI pada plant kecepatan dengan saturasi $|u| \le 1{,}2$, dengan dan tanpa \textit{anti-windup}.
	\item Mengapa penguatan derivatif efektif membesar bila $T$ dikecilkan, dan bagaimana filter derivatif mengatasinya?
\end{enumerate}
